%%%%%%%%%%  scan idx 475  |  folio 462  %%%%%%%%%%
\SGAPrintPage{462}%
(4.8.1), we obtain:
\[
  f_*(j_*(y'))+f_*(f^*(x))\in\Fil^{k+1}K^{\bullet}(X)_{\mathbb Q},
\]
or, by 3.6:
\[
  i_*(y)+x\in\Fil^{k+1}K^{\bullet}(X)_{\mathbb Q}.
\]

To prove relation ii), we must show that:
\[
  y'-g^*(g_*(y'))\lambda_{-1}(F)\in\Fil^kK^{\bullet}(Y')_{\mathbb Q}.
\]
Put \(z=y'-g^*(g_*(y'))\lambda_{-1}(F)\). We have:
\[
\begin{aligned}
j_*(z)
  &=j_*(y')-j_*\bigl(g^*(g_*(y'))\lambda_{-1}(F)\bigr)\\
  &=j_*(y')-f^*(i_*(y))\qquad\text{by 3.4.}
\end{aligned}
\]
Now relation i) gives:
\[
  f^*(i_*(y))+f^*(x)\in\Fil^{k+1}K^{\bullet}(X')_{\mathbb Q}.
\]
Taking (4.8.1) into account, we therefore obtain:
\[
  j_*(z)\in\Fil^{k+1}K^{\bullet}(X')_{\mathbb Q},
\]
whence:
\[
  j^*(j_*(z))=z(1-L)\in\Fil^{k+1}K^{\bullet}(Y')_{\mathbb Q}.
\]

Let \(\bar z\) be the image of \(z\) in \(\operatorname{Gr}^{k-1}K^{\bullet}(Y')_{\mathbb Q}\), and \(\overline{(1-L)}\) that of \((1-L)\) in \(\operatorname{Gr}^{1}K^{\bullet}(Y')_{\mathbb Q}\). We therefore have: \(\overline{z(1-L)}=0\). We can then apply (VI 5.11) (which remains valid after tensoring with \(\mathbb Q\), since its proof uses only the structure of \(K^{\bullet}(Y')\) as a free \(K^{\bullet}(Y)\)-algebra), and hence find \(z'\in K^{\bullet}(Y)_{\mathbb Q}\) such that:
\[
  z-g^*(z')\lambda_{-1}(F)\in\Fil^kK^{\bullet}(Y')_{\mathbb Q}.
\]
Applying \(g_*\), and using \(g_*(z)=0\), we find:
\[
  z'\in\Fil^{k+1-d}K^{\bullet}(Y)_{\mathbb Q},
\]
whence, since \(\lambda_{-1}(F)\in\Fil^{d-1}K^{\bullet}(Y')_{\mathbb Q}\):
\[
  z\in\Fil^kK^{\bullet}(Y')_{\mathbb Q},
\]
which completes the proof.

%%%%%%%%%%  scan idx 476  |  folio 463  %%%%%%%%%%
\SGAPrintPage{463}%
Suppose now that \(X\) is noetherian and \(Y\) regular, which implies that \(v\) is surjective (3.7). We saw in the proof of 3.7 that for every \(x'\in K^{\bullet}(X')\) there is \(y'\in K^{\bullet}(Y')\) such that:
\[
  x'=f^*(f_*(x'))+j_*(y').
\]
The same property clearly remains true after tensoring with \(\mathbb Q\). Consequently, for \(x'\in\Fil^kK^{\bullet}(X')_{\mathbb Q}\) there is \((y',x)\in K^{\bullet}(Y')_{\mathbb Q}\times K^{\bullet}(X)_{\mathbb Q}\) such that:
\[
  x'=v(y',x),\qquad x\in\Fil^kK^{\bullet}(X)_{\mathbb Q}.
\]
Indeed, it suffices to take \(x=f_*(x')\).

Filter the ring \(K^{\bullet}(Y')_{\mathbb Q}\times K^{\bullet}(X)_{\mathbb Q}\) defined in 4.7 by putting:
\[
  \Fil^k\bigl(K^{\bullet}(Y')_{\mathbb Q}\times K^{\bullet}(X)_{\mathbb Q}\bigr)
  =
  \Fil^{k-1}K^{\bullet}(Y')_{\mathbb Q}\times \Fil^kK^{\bullet}(X)_{\mathbb Q}.
\]
It follows immediately from (4.7.1) that this is a ring filtration. Thus, to show that for every \(k\) the homomorphism:
\[
  \Fil^k\bigl(K^{\bullet}(Y')_{\mathbb Q}\times K^{\bullet}(X)_{\mathbb Q}\bigr)
  \xrightarrow{v}
  \Fil^kK^{\bullet}(X')_{\mathbb Q}
\]
is surjective, it suffices to show that for \(x'= \gamma^k(x'_1)\), with \(\epsilon(x'_1)=0\), there is \((y',x)\in\Fil^k(K^{\bullet}(Y')_{\mathbb Q}\times K^{\bullet}(X)_{\mathbb Q})\) such that \(x'=v(y',x)\).

By the preceding argument, there is \((y'_1,x_1)\in K^{\bullet}(Y')_{\mathbb Q}\times K^{\bullet}(X)_{\mathbb Q}\) such that \(x'_1=v(y'_1,x_1)\) and \(x_1\in\Fil^1K^{\bullet}(X)_{\mathbb Q}\). By 4.7, we then have:
\[
  x'=v\bigl(\gamma^k_*(y'_1,x_1)\bigr).
\]
Now:
\[
  \gamma^k_*(y'_1,x_1)
  =
  \left(\sum_{p=1}^{k}\gamma^{k-p}(x_1)\gamma^p(L,y'_1),\;
  \gamma^k(x_1)\right)
\]
(an immediate consequence of (4.7.2)). Since \(x_1\) has zero augmentation, \(\gamma^k(x_1)\in\Fil^kK^{\bullet}(X)_{\mathbb Q}\), and, using V 5.7,
\[
  \sum_{p=1}^{k}\gamma^{k-p}(x_1)\gamma^p(L,y'_1)
  \in\Fil^{k-1}K^{\bullet}(Y')_{\mathbb Q},
\]
which gives the result.

%%%%%%%%%%  scan idx 477  |  folio 464  %%%%%%%%%%
\SGAPrintPage{464}%
\sgaStableTarget{sga6.e07.rem.4-9}\textbf{Remark 4.9.} When \(X\) is noetherian and \(Y\) regular, proceeding as in 4.7 we can describe the graded ring structure of \(\operatorname{Gr}^*K^{\bullet}(X')_{\mathbb Q}\) in terms of the graded rings \(\operatorname{Gr}^*K^{\bullet}(X)_{\mathbb Q}\), \(\operatorname{Gr}^*K^{\bullet}(Y)_{\mathbb Q}\), the additive homomorphisms \(i^{\gr}\) and \(i_{\gr}\) induced by \(i^*\) and \(i_*\), and the element \(c^{\bullet}(N)\) of \(\operatorname{Gr}^*K^{\bullet}(Y)\). Likewise, the same data describe the homomorphisms \(f^{\gr},f_{\gr},j^{\gr},j_{\gr}\) defined from \(f^*,f_*,j^*,j_*\) by passage to the associated gradeds.

\sgaStableTarget{sga6.e07.sec.4.10}\subsection*{4.10.}

Let \(X\) be a \emph{regular} scheme having an ample sheaf. We know (IV\qquad) that the Grothendieck groups formed with locally free \(\underline{O}_X\)-Modules of finite type and with coherent \(\underline{O}_X\)-Modules are isomorphic. We use this fact to define a ``topological'' filtration on \(K^{\bullet}(X)\), taking \(\Fil^i_{\mathrm{top}}K^{\bullet}(X)\) to be the subgroup generated by the classes of coherent sheaves whose support has codimension \(\ge i\); it is easy to show that \(\Fil^i_{\mathrm{top}}K^{\bullet}(X)\) is generated by the classes of the sheaves \(\underline{O}_Z\), where \(Z\) ranges over the integral closed subschemes of \(X\) of codimension \(\ge i\)(X\qquad).

\sgaStableTarget{sga6.e07.prop.4-11}\textbf{Proposition 4.11.} \uline{Let} \(X\) \uline{be a regular scheme having an ample sheaf}. \uline{Then} \uline{for every} \(k\ge0\):
\[
  \Fil_{\mathrm{top}}^kK^{\bullet}(X)_{\mathbb Q}
  \subset
  \Fil_\lambda^kK^{\bullet}(X)_{\mathbb Q}.
\]
\((\mathrm{N.B.}\) The reverse inclusion, valid without tensoring with \(\mathbb Q\), will be proved in X\qquad).

By 4.10, it suffices to show that for an integral closed subscheme \(Z\) of \(X\), of codimension \(k\), we have:
\[
  \operatorname{cl}(\underline{O}_Z)\in\Fil^kK^{\bullet}(X)_{\mathbb Q}.
\]

For this, we proceed by noetherian induction on the integral closed subschemes of \(X\). Suppose therefore that the property is proved for the integral closed subschemes of \(Z\), and let us prove it for \(Z\). Let \(\underline{J}\) be the Ideal of \(\underline{O}_X\)

%%%%%%%%%%  scan idx 478  |  folio 465  %%%%%%%%%%
\SGAPrintPage{465}%
defining \(Z\), and let \(z\) be the generic point of \(Z\). Since \(Z\) is integral, \(\underline{J}_z\) is the maximal ideal of \(\underline{O}_{X,z}\), and since \(X\) is regular, \(\underline{J}_z\) can be generated by an \(\underline{O}_{X,z}\)-regular sequence. We can therefore find a sequence of sections of \(\underline{J}\) in a neighborhood of \(z\) forming an \(\underline{O}_X\)-regular sequence. Consequently, there is an open neighborhood \(U\) of \(z\) such that \(Z'=X-U\) is contained in \(Z\), and \(\underline{J}|U\) is a regular ideal of \(\underline{O}_U\). The immersion \(Z\cap U\xrightarrow{i}U\) is then a regular immersion, with codimension \(\dim(\underline{O}_{X,z})=\codim(Z,X)\) (EGA IV 5.1.2). Furthermore, we have the exact sequence (IV\qquad):
\[
  K^{\bullet}(Z')\xrightarrow{j_*}K^{\bullet}(X)
  \xrightarrow{u^*}K^{\bullet}(U)\to0,
\]
where \(j\) is the immersion \(Z'\to X\) and \(u\) the immersion \(U\to X\); tensoring with \(\mathbb Q\) gives an analogous sequence. Now \(u^*(\underline{O}_Z)=i_*(\underline{O}_{Z\cap U})\), so by 4.6 the image \(u^*(1_Z)\) of the class of \(\underline{O}_Z\) under \(u^*\) belongs to \(\Fil^kK^{\bullet}(U)_{\mathbb Q}\). Since \(u^*\) is a surjective \(\lambda\)-homomorphism, it induces a \emph{surjection} \(\Fil^kK^{\bullet}(X)\longrightarrow \Fil^kK^{\bullet}(U)\), so \(u^*(1_Z)=u^*(x)\) with \(x\in\Fil^kK^{\bullet}(X)\); there is therefore \(y\in j_*(K^{\bullet}(Z'))\) such that:
\[
  \operatorname{cl}(\underline{O}_Z)-y\in\Fil^kK^{\bullet}(X)_{\mathbb Q}.
\]
Now \(Z'\subset Z\), and since \(z\notin Z'\), \(\codim(Z',X)>k\). Consequently, \(y\) can be expressed as a combination of classes of sheaves \(\underline{O}_{Z''_\alpha}\), where the \(Z''_\alpha\) are integral closed subschemes of \(Z\), distinct from \(Z\). By the induction hypothesis, \(y\in\Fil^{k+1}K^{\bullet}(X)_{\mathbb Q}\), which completes the proof.

\sgaStableTarget{sga6.e07.bibliography}\subsection*{BIBLIOGRAPHY}
\begin{enumerate}[label={[\arabic*]}]
\item Cartan--Eilenberg, Homological Algebra.
\end{enumerate}



% Expose VIII macro additions
\providecommand{\etale}{\'etale}
\providecommand{\Ch}{\operatorname{Ch}}


